Tutte le fonti consultate per la stesura del presente documento sono riportate nella bibliografia. Gli strumenti di intelligenza artificiale generativa (quali \textit{OpenAI ChatGPT} e \textit{Google Gemini}) sono stati impiegati esclusivamente per la revisione stilistica e il controllo ortografico e grammaticale.\newline

\noindent Ad esempio il testo originale stilato dal gruppo:\newline

\noindent \textit{Il mondo dell'agricoltura è sempre stato interessato da problematiche di ottimizzazione delle risorse uno degli esempi più critici è proprio quello dell'utilizzo di acqua; si pensi a momenti di grande secchezza e quindi al razionamento delle risorse idriche nei periodi più caldi dell'anno. Esistono anche problemi riguardanti all'uso di fertilizzanti e dei pesticidi entrambi da regolamentare per sottostare a normative statali o sovrastatali (ad esempio leggi comunitarie europee)}.\newline

\noindent E' stato revisionato in:\newline

\noindent\textit{Il settore agricolo è storicamente caratterizzato dalla necessità di ottimizzare le risorse, tra cui la gestione idrica rappresenta una delle criticità maggiori. In contesti di siccità prolungata, il razionamento dell'acqua durante i mesi estivi diventa una sfida determinante per la sopravvivenza delle colture. Parallelamente, l'impiego di fertilizzanti e pesticidi richiede un monitoraggio rigoroso, non solo per ragioni ambientali ma anche per garantire la conformità alle normative nazionali e comunitarie (UE), sempre più restrittive in termini di sostenibilità}.
